% Official contribution: https://conferences.iaea.org/event/498/contributions/43621/
% Scheduled date: 2026-10-27.
% Official description is "TBD"; treated as no abstract.
\paragraph{Abstract}
% No substantive abstract published at the time of preparation.
This presentation provides an overview of recent developments at the NEA Data Bank for 3 different topics around international collaborations and project management.

First topic are the conceptual requirements for a collaboration in the 21st century. Key requirements identified are transparency, and tracability of contributions, persistence of releases as well as common project management practices and data structures. Examples are provided from the recent NEA SCI project modernisation.

Secondly, the need for open data standards and facilitated data contribution to international projects has been identified as key requirements. Recent implementations on this from the NEA are an Exfor Submission Form (more details in \ref{subsec:EXFOR_Streamlined}).  

As third topic, persistent identifiers and integration into field-specific resource networks are identified as essential to make resources findable and accessible via standard interfaces. Here, the NEA DB launched an InvenioRDM instance which integrates records into the DataCite framework, providing DOIs and interlinking with related resources.

Finally, the presentation triggered discussions on community practices, sustainability, project ownership and future developments paths of legacy tools and datasets and new developments.

\mbox{}\par

\paragraph{Discussion}
% Source: Webex-Summary/day2.pdf; use the meeting date, not the filename.
% Also checked against published slides; do not attribute later JADE/git-annex discussion here.
Julia Sprenger's contribution described NEA initiatives for tracking contributions, community ownership and persistent releases~\cite{neaSlides}. Modular schemas and automated conversion, validation and documentation were presented as building blocks shared between projects. Discussion concerned making collaboration through Git easier for users unfamiliar with these tools and connecting repositories, APIs and user-facing services. The EXFOR form illustrated a contribution pathway still requiring decisions on authentication and backend integration. The public slides also acknowledge Anastasia Georgiadou.
